\chapter{Literature Review}

The rapid advancement of artificial intelligence and computer vision technologies has significantly transformed modern home appliances and laundry systems. Manufacturers are increasingly adopting smart technologies to improve washing efficiency, enhance fabric care, reduce resource consumption, and provide personalized laundry experiences. Smart washing machines, fabric classification systems, automated detergent dispensing mechanisms, and IoT-enabled home appliances play a major role in improving laundry quality, sustainability, and overall user satisfaction. This chapter reviews recent literature related to smart home appliances, fabric detection technologies, automated resource management, and human-machine interaction in laundry systems.

\noindent
The technologies used in smart washing machines include computer vision systems, deep learning algorithms for fabric classification, IoT-based monitoring and control, sensor integration for load detection, and automated dispensing mechanisms. However, the literature also identifies limitations such as high implementation costs, accuracy challenges in fabric identification, technical complexity in real-world conditions, data privacy concerns with connected devices, and the need for extensive training datasets. In some cases, over-automation may reduce user control and customization options.

\section{Smart Home Appliances and Laundry Technology}

This section analyzes how smart home appliances, particularly smart washing machines, influence both user experience and laundry outcomes. The functionality of smart laundry systems includes automated cycle selection, remote monitoring and control through mobile applications, energy and water consumption tracking, maintenance alerts, and integration with home automation ecosystems. These systems aim to improve convenience, reduce resource waste, enhance fabric care, and provide users with greater control over their laundry processes. The literature emphasizes improved user convenience, optimized resource usage, and enhanced garment longevity as key outcomes of smart washing machine implementation.

\noindent
Technologies used in smart laundry systems include IoT connectivity, embedded sensors, mobile applications, cloud-based data storage, and machine learning algorithms for pattern recognition. Research emphasizes convenience, energy efficiency, and personalized washing as key benefits of smart appliance adoption. However, limitations include high initial costs, compatibility issues with existing home infrastructure, user learning curves, and concerns about long-term reliability of smart features. In some cases, complex interfaces may reduce accessibility for elderly users.

\section{Fabric Detection and Computer Vision in Laundry Applications}

This research focuses on applying computer vision and deep learning techniques for automatic fabric identification in laundry systems. The functionality of fabric detection systems includes image capture of clothing items, preprocessing of visual data, feature extraction (texture, color, pattern), and classification of fabric types using trained neural networks. These systems aim to eliminate manual sorting, reduce dependency on fabric labels, and enable customized wash cycles based on detected materials.

\noindent
The technologies used include high-resolution cameras, image processing libraries (OpenCV), convolutional neural networks (CNNs) for image classification, transfer learning with pre-trained models (ResNet, VGG, MobileNet), and embedded vision systems for real-time processing. Despite significant progress, the literature identifies limitations such as difficulty distinguishing similar fabric textures (e.g., different cotton blends), challenges with varying lighting conditions inside the washing machine drum, occlusion when clothes overlap, and the need for large, diverse fabric datasets. Additionally, fabric detection accuracy may decrease with wet or wrinkled clothing.

\section{Automated Resource Management Systems}

This scoping review examines various technologies for automated detergent and water dispensing in washing machines. The functionality of these systems includes load weight sensing, soil level detection, fabric type consideration, and precise dispensing of detergent and water based on calculated requirements. These mechanisms aim to minimize resource wastage, prevent detergent residue on clothes, reduce environmental impact, and ensure optimal cleaning results.

\noindent
The technologies discussed in the review include load cells and weight sensors for load detection, turbidity sensors for soil level measurement, conductivity sensors for detergent concentration monitoring, peristaltic pumps for precise liquid dispensing, and microcontroller-based control systems. However, the review highlights limitations such as sensor calibration challenges, variability in detergent viscosity and concentration, difficulty measuring liquid detergents accurately, and increased mechanical complexity. Furthermore, existing systems often rely on generalized algorithms rather than fabric-specific adjustments, leading to suboptimal resource usage.

\section{AI and Machine Learning in Appliance Optimization}

This study explores how machine learning algorithms optimize washing machine performance and user experience. The functionality of AI-optimized laundry systems includes learning from user preferences, adapting to typical laundry patterns, predicting maintenance needs, optimizing cycle parameters based on historical data, and providing personalized recommendations. These intelligent systems enhance efficiency, reduce energy consumption, and improve overall user satisfaction through adaptive behavior.

\noindent
The technologies used include supervised and unsupervised learning algorithms, reinforcement learning for cycle optimization, predictive maintenance models, user behavior analytics, and cloud-based machine learning pipelines. However, limitations include the requirement for substantial user data to train effective models, privacy concerns with data collection, computational resource constraints on embedded devices, and the "black box" nature of some AI decisions that may reduce user trust. In some cases, over-personalization may lead to unexpected behavior that confuses users.

\section{Summary}

The proposed AI-Integrated Washing Machine aims to overcome these limitations by providing:

\begin{itemize}
    \item A vision-based fabric detection system that automatically identifies clothing materials without manual input
    \item Intelligent resource calculation and dispensing based on detected fabric types, load weight, and soil level
    \item Customized wash cycle generation for optimal fabric care and cleaning results
    \item A user-friendly interface displaying detected fabrics and recommended settings
    \item Integration of weight sensing and soil detection for comprehensive laundry optimization
    \item Practical implementation suitable for residential and small commercial applications
\end{itemize}

\noindent
Unlike many complex smart laundry systems that require extensive user configuration or rely on generalized cycles, the AI-Integrated Washing Machine focuses on automation, precision, and fabric-specific care. By combining computer vision, machine learning, and automated dispensing with an intuitive interface, the system enhances garment longevity while optimizing resource consumption and improving user convenience.

\noindent
Thus, this project contributes to the growing field of intelligent home appliances by offering a practical, accessible, and fabric-centered laundry solution.